\phantomsection
\addcontentsline{toc}{subsection}{Canales de Distribución (CD)}
\subsection*{Canales de Distribución (CD)}

Es la forma en la que especificamos como será el contacto con los clientes, para ello, el modelo de distribución de nuestro proyecto será principalmente a través de un sitio web intuitivo. Junto a divulgación en redes sociales (haciendo presencia en plataformas donde se concentra el público objetivo como Facebook, TikTok, etc.) y marketing de contenidos para captación. 

Además, dentro de las posibilidades se contemplan alianzas con universidades, ONGs y fintech locales para amplificar nuestra difusión (mediante talleres o seminarios). Ejemplos son la Secretaría Distrital del Hábitat el cuál promueve cursos de educación financiera \parencite{SecretariaBogota01}, el Banco de la República que cuenta con programas de educación financiera \parencite{BanRepEduca01}, y la Universidad del Rosario que posee programas similares para los jóvenes \parencite{UniverRosario01}. Por último, publicidad en línea paga (mediante Google Ads).